\section{Exact phase selection using the free onset}
\label{sec:effective-phase}
The ceiling in the all-time construction is a literal integer, not a
safe upper integer. Approximation alone cannot decide it at an arbitrary
prescribed onset if equality has not been excluded. Note~60 \cite{n60}
uses the existing freedom to enlarge that onset before choosing $X_0$
or the source. The unwrapped phase here is fixed by
\begin{align}
 w&=(1+ix)/2,&p&=w+32\alpha(p),&v&=p-w,&B&=1-32\alpha'(p),\notag\\
 \log S&=\log M(p)-v^2/64-\tfrac12\Log B,&
 \Theta(64,x)&=\Im\log S.
 \label{eq:pi-normalization}
\end{align}
The log-gamma branch is analytic in the right half-plane, real on the
positive axis and normalized by $\log\Gamma(1)=0$. No principal
argument of the final small arm replaces it.

\begin{lemma}[Full-saddle phase remainder]\label{lem:pi-phase}
For $x\ge e^{1000}$, $\lambda=\log x$ and $L=\log(x/(4\pi))$,
\begin{equation}
 \Theta(64,x)=\frac x4(L-1)+4\pi L+\frac{7\pi}8+E(x),\qquad
 |E(x)|\le2^{16}(1+\lambda)^2/x.
 \label{eq:pi-expansion}
\end{equation}
\end{lemma}
\begin{proof}
The proved saddle rectangle gives $\Re p>2$, $|p|\ge x/2$,
$|v|\le32\lambda$,
$|\alpha(p)-\frac12\Log(p/(2\pi))|\le4/x$ and
$|B-1|\le97/x$. Stirling on that right-half-plane branch gives
\[
 \log M(p)=F(p)+\tfrac32\Log p+\log(\sqrt\pi/8)
                       +\Log(1-1/p)+R_1(p/2),
 \quad F(z)=\tfrac z2[\Log(z/(2\pi))-1].
\]
Both last errors have modulus at most $4/x$: use the convergent
logarithm series and the sector remainder
$\sec^2(\arg z/2)/(12|z|)$ \cite{DLMFGammaOuter2026}.
The identity first holds for positive $p$ and then by continuation;
there is no undetermined phase translation.

Put $q=ix/2$, $h=p-q=v+1/2$ and
$\alpha_0=\frac12\Log(q/(2\pi))=L/2+i\pi/4$. The segment from
$q$ to $p$ avoids the cut, has $|z|\ge x/4$ and $|h|\le64\lambda$.
Taylor's integral formula, with $F''=1/(2z)$, gives
\[
 |F(p)-F(q)-\alpha_0h|\le4096\lambda^2/x,\quad
 |\Log p-\Log q|\le256\lambda/x,
\]
and comparison of $\alpha$ at $p$ gives
$|\alpha(p)-\alpha_0|\le129\lambda/x$,
$|v-32\alpha_0|\le2^{13}\lambda/x$. Complete the square exactly:
\[
 \alpha_0v-v^2/64=16\alpha_0^2-(v-32\alpha_0)^2/64.
\]
The last error is at most $\lambda^2/x$, and
$|\Log B|/2\le97/x$. Thus
\[
 \log S=F(q)+\alpha_0/2+\tfrac32\Log q+16\alpha_0^2
                          +\log(\sqrt\pi/8)+R,
 \quad |R|<5000\lambda^2/x.
\]
The respective imaginary parts are $x(L-1)/4$, $\pi/8$,
$3\pi/4$, and $4\pi L$. This proves the claimed expansion of the
implicit full saddle, including its displacement payment.
\end{proof}

\begin{theorem}[Literal phase index with free onset]
\label{thm:phase-selector}
Let $K\ge2048$ be a certified integer upper bound for every onset
exponent gate of Theorem~\ref{thm:all-time}, including the enlarged
positive band and pre-switch bridge gates. Among $U_j=2^{K+j}$,
$0\le j\le3$, at least one has phase distance greater than $1/100$
from an integer. A finite selector returns a possibly different onset with
\[
 \operatorname{dist}(\Theta(64,U_j)/\pi,\mathbb Z)>1/200
\]
and its exact $n=\lceil\Theta(64,U_j)/\pi\rceil$.
For $P_*=K+3+\lceil\log_2(K+4)\rceil+20$, sufficient selector
cost and storage bounds are $2^{100}(P_*+1)^{10}$ bit operations
and $2^{20}(P_*+1)^3$ bits. They exclude determining $K$ and the
remaining construction.
\end{theorem}
\begin{proof}
Writing $Q_l=\Theta(64,2^l)/\pi$, Lemma~\ref{lem:pi-phase} gives
\begin{align*}
 Q_l&=A_l+r_l,\\
 A_l&=\frac{2^{l-2}}\pi[(l-2)\log2-\log\pi-1]
                 +4[(l-2)\log2-\log\pi]+7/8,\\
 |r_l|&\le2^{16}(1+l\log2)^2 2^{-l}.
\end{align*}
For $l\ge K$, the error is $<2^{-26}$: use $\lambda>K/2\ge1024$,
$e^\lambda\ge\lambda^8/8!$ and $8!<2^{16}$.
For $c=(-4,8,-5,1)$,
$\sum c_j2^j=\sum c_jj2^j=\sum c_j=0$, $\sum c_jj=1$,
and $\sum|c_j|=18$. Hence
\[
 D=Q_{K+3}-5Q_{K+2}+8Q_{K+1}-4Q_K,
 \qquad |D-4\log2|<1/100.
\]
The positive atanh series proves $2/3<\log2<7/10$; consequently
$\operatorname{dist}(D,\mathbb Z)>19/100$. Nearest integers to
the four $Q$'s would give
$\operatorname{dist}(D,\mathbb Z)\le18\max_j
\operatorname{dist}(Q_{K+j},\mathbb Z)$, so one distance exceeds
$19/1800>1/100$.

Here is a terminating rational selection, which may reject equalities.
For each $l=K,\ldots,K+3$, set
$P=l+\lceil\log_2(l+1)\rceil+20$, $N=P+10$. Compute
\begin{align*}
 \pi_N&=16\sum_{r<N}\frac{(-1)^r}{(2r+1)5^{2r+1}}
         -4\sum_{r<N}\frac{(-1)^r}{(2r+1)239^{2r+1}},\\
 \widehat a&=2\sum_{r<N}\frac{3^{-2r-1}}{2r+1}.
\end{align*}
Round $\pi_N$ to grid $2^{-P-4}$, obtaining $3\le\widehat\pi\le4$,
$|\widehat\pi-\pi|<2^{-P-4}$; the alternating Machin tails are
$<20\cdot25^{-N}<2^{-P-6}$. Set
$\rho=(\widehat\pi-1)/(\widehat\pi+1)\in[1/2,3/5]$ and
$\widehat b=2\sum_{r<N}\rho^{2r+1}/(2r+1)$. The atanh tails
give $|\widehat a-\log2|,|\widehat b-\log\pi|<2^{-P}$,
including the separate logarithm-input error
$|\widehat\pi-\pi|/3$.

Form exact rationals
\[
 \widehat L=(l-2)\widehat a-\widehat b,\quad
 \widehat Q=2^{l-2}(\widehat L-1)/\widehat\pi+4\widehat L+7/8,
 \quad I_l=[\widehat Q-1/1000,\widehat Q+1/1000].
\]
The leading-phase constant error obeys
$|\widehat Q-A_l|\le2^ll2^{-P}<2^{-20}$; adding $r_l$ gives
$|\widehat Q-Q_l|<2^{-19}$. Thus $I_l$ encloses $Q_l$.
Compute $n=\lfloor\inf I_l\rfloor+1$ by rational division and
accept the first interval with
\[
 n-1+1/200<\inf I_l,\qquad \sup I_l<n-1/200.
\]
A rational tie means rejection. The separated member has whole
interval distance $>1/100-2^{-19}-1/1000>1/200$, so one test passes.
Its two inequalities prove the literal ceiling and strict margin.

Common denominators $(2N)!5^{2N}239^{2N}$ and $(2N)!3^{2N}$
suffice for Machin and $\log2$. If $\rho=R/S$, its numerator and
denominator have at most $P+8$ bits after the prescribed rounding;
$(2N)!S^{2N}$ suffices for $\log\widehat\pi$. All operands,
scaled sums and comparisons fit $2^{10}(P_*+1)^2$ bits, using
fewer than $2^{10}(P_*+1)^2$ integer operations. Schoolbook
arithmetic and streaming give the stated generous cost/storage bounds.
The integer returned has $O(K+\log(K+1))$ bits.
\end{proof}

The order of choices is part of the theorem. First fix $\eta,\theta$,
$C_+,T_+$, then $b=T_++1$, $\delta=\theta/64$, $T_-$,
$K_{\rm mid}$ and $K_{\rm br}$. With the original notation of the
compact table, $K$ must dominate
\begin{align*}
 16384+\lceil\log_2(4R_b^*+4D_{\rm mov}+1024)\rceil
           +\lceil8\log_2(2+C_b^+/\delta)\rceil,\\
 16384+\lceil8\log_2(2+2^{64}K_{\rm br}/\delta)\rceil,
 \qquad2048,
\end{align*}
where $C_b^+=(1+c_b)C_b$, and any auxiliary gate remains included.
Their permitted strict upper integers require no exact-real ceiling.
All four candidate exponents retain the monotone gates, including
$a_X\ge2^{k-1}$, $L\le k+3$, $C_b^+\le\delta2^{k/8}$ and
$(k+3)^8\le2^{k/2}$. Select the winner before fixing $X_0$; then
invert $\Theta(64,X_0)=n\pi$ in $[U_0,U_0+32]$, and choose
the moving gap and one source from the resulting full minimum.
No source chosen at a rejected onset is reused. This resolves the
discrete input for the authorized free-onset family. Equality at an
arbitrary already prescribed dyadic exponent remains outside the result.

\section{A complete evaluator and continuous-filter ledger}
\label{sec:complete-evaluation}
Notes~59--61 \cite{n59,n60,n61} give analytical algorithms for the
same complete tuple \eqref{eq:ua-contract}. A certified descriptor
contains the exact dyadic source, $U,m,\beta$, the exact phase index
from Section~\ref{sec:effective-phase}, encoded $\eta>0$ with a positive
lower bound, thresholds, widths and branch/guard margins. Independent
scalar data are rational/dyadic or come with computable approximation
routines and separately charged costs. Unspecified exact real inputs
are not algorithmic data. No coefficients or research numerical runs
are produced by the following proofs.

\subsection{Cancellation-safe evaluation of the stored source}
The expanded source satisfies
\begin{equation}
 D_f=\sum_k|d_k|\le16\cdot49^m<2^{6m+4},\qquad
 0\le f(u)<100\quad(0\le u\le U).
 \label{eq:ev-source-caps}
\end{equation}
These are different caps: the first pays coefficient cancellation;
the second supplies better coordinate and smoothness bounds.
For $a=tU^2$, $c_\pm=\pm izU$, absolute convergence gives
\begin{equation}
 \partial_t^h\partial_z^jG=\frac{U^{2h+j+1}}2\sum_kd_k
 [i^jM_{k+2h+j}(a,c_+)+(-i)^jM_{k+2h+j}(a,c_-)].
 \label{eq:ev-jets}
\end{equation}
There is no division by $a$ or $c$ in the series
\eqref{eq:ac-entire-moments}, so this evaluates the same field across
$t=0,z=0$ without assuming a stable forward recurrence.

\begin{theorem}[Finite working-bit source evaluator]\label{thm:ev-source}
Fix $2h+j\le J_e$, with $J_e=8$ sufficient below, and a rational
query/slack rectangle $-\tau\le t\le b$, $|x|\le X$, $|y|\le Y\le4$.
For $0<\nu\le1$, set $s_b=\max(0,b)$ and
\begin{align}
 A_e&=(\tau+s_b+2)U^2+U(X+Y+2),&T_e&=\lceil A_e\rceil,\notag\\
 p_e&=\left\lceil\log_2(4\cdot2^{6m+4}U^{J_e+1}/\nu)\right\rceil,
 &D_e&=8(T_e+1)+p_e-1.
 \label{eq:ev-truncation}
\end{align}
Truncate each moment series at $r+s\le D_e$. There is a finite
exact-rational implementation with total jet error $<\nu$, including
coordinate rounding and final output rounding. With the explicit
integer $K_e$ defined below, sufficient operation, operand and
schoolbook bit bounds are respectively
$2^{20}K_e^5$, $2^{20}K_e^3$ bits and $2^{100}K_e^{11}$.
\end{theorem}
\begin{proof}
With $N=D_e+1\ge8A_e$, the joint exponential tail is
\[
 e^{A_e}A_e^N/N!\le2^{2A_e-N}\le2^{-p_e},
\]
using $N!\ge(N/e)^N$ and $e<3$. Multiplication by
$D_fU^{J_e+1}$ pays $\nu/4$, including all coefficient cancellation.
For rounding the three coordinates, every next real derivative is at
most
\[
 V_J=100U^{J_e+3}e^{(s_b+1)U^2+(Y+1)U}.
\]
Choose common coordinate denominator $2^{B_e}$ with
$B_e\ge\max(0,\lceil\log_2(12V_J/\nu)\rceil)$;
the displacement segment is in the declared slack box and costs
$\nu/4$. Exact dyadic coordinates need no displacement debit.

Set $H_e=2m+J_e+2D_e+2$ and
\[
 Q_e=2^{2\beta+B_eD_e+1}D_e!H_e!,\qquad
 p_{\rm out}=\lceil\log_2(8/\nu)\rceil.
\]
Each series term has this common denominator, since
$D_e!/(r!s!)$ and $H_e!/(k+2h+j+2r+s+1)$ are integers.
Generate Gaussian-integer powers, sum exactly, and round only the final
complex output to $p_{\rm out}$ fractional bits, costing $\nu/4$.
Put
\begin{align*}
 \ell_e&=\lceil\log_2(U+T_e+D_e+H_e+m+\beta+B_e+p_{\rm out}+J_e+8)\rceil,\\
 L_e&=64[2\beta+B_e(D_e+1)+(D_e+H_e+J_e+4)\ell_e
                      +6m+2T_e+p_{\rm out}+16],\\
 K_e&=2+U+m+\beta+B_e+D_e+J_e+T_e+p_{\rm out}.
\end{align*}
The denominator length is at most
$2\beta+B_eD_e+1+D_e\log_2D_e+H_e\log_2H_e$;
numerator powers cost at most
$D_e(B_e+\log_2(T_e+1)+3)$ bits. Absolute partial sums add at most
$\log_2D_f+2T_e+(J_e+1)\log_2U+4$ bits. Thus all operands fit
$L_e\le2^{20}K_e^3$. There are $O(mD_e^2)$ terms; even naive
power/factorial scaling fits $2^{20}K_e^5$ operations. Schoolbook
arithmetic gives the stated bit bound. If query encoding length is
$I$, add at most $2^{20}(I+B_e+1)^2$ reading/rounding work and
$I+B_e$ storage. Large rational denominators are paid, not dropped.
\end{proof}
These are exact-field evaluation bounds. Source-to-$H$ accuracy is
still only \eqref{eq:ac-four-jets}. The source is active at every
real height on the negative attachment, middle interval and initial
switch; its active time is bounded. No spatial cutoff is introduced.

\subsection{All nonsource leaves and the finite jet recursion}
For the complex saddle retain its paid contraction rectangle and
$|B_s|\ge1/2$. Inexact fixed-point iteration with local error
$\eta_{\rm it}$ obeys
\[
 |p_k-p|\le2^{-k}D_{\rm rect}+2\eta_{\rm it}.
\]
Choosing $k\ge\lceil\log_2(4D_{\rm rect}/\delta)\rceil$ and
$\eta_{\rm it}\le\delta/8$ pays a $\delta$ value target, with
separate final rounding. Projection onto a closed rational rectangle
inside the inherited strict margins preserves its domain. The actual
arm logarithm is
\begin{equation}
 \log S=\Log p+\Log(p-1)-\tfrac p2\log\pi+\log\Gamma(p/2)
       -\log16-(p-w)^2/t-\tfrac12\Log B_s.
 \label{eq:ev-arm-log}
\end{equation}
It is used only at $t\ge63$. In the truncated real Taylor ring in
$t,x,y$, each highest positive-degree coefficient of
$p-w-t\alpha(p)/2=0$ has multiplier $B_s$; solve with that paid
floor. Product, inverse and exponential recurrences then give all
fixed jets. This differentiates the actual implicit equation.

With $n$ fixed, phase inversion is also constructive. On a width-$w$
bracket with $\Theta_x\ge c>0$, evaluate the midpoint residual to
error $cw/16$. A certified sign chooses a half-bracket. Otherwise its
true residual has modulus $\le cw/8$, so the interval of radius
$w/8$ about the midpoint contains the root. Each step halves width
without a root-equality test. The velocity is
$X'=-\Theta_t/\Theta_x$ and higher finite jets use the same
implicit-coefficient rule. This does not select $n$; that is paid by
Theorem~\ref{thm:phase-selector}.

The physical source saddle on $[63,B_q]$ uses
\[
 L_s=\tfrac14\log[(2tL_s+9)/(4\pi)],\quad
 1/4\le L_s\le B_q+1,\quad
 \rho_q=B_q/(B_q+18)<1.
\]
Projected inexact iteration takes at most
$\lceil(B_q+18)\log(4(B_q+1)/\delta)/18\rceil$ steps when its
local error is at most $(1-\rho_q)\delta/4$. Compute
$\kappa=8tL_s-2t+36\ge36$, $\log M_0$, and the actual
$r=a(t)-\beta_q(t)x^2$ with their finite jets. At any late query
take $B_q\ge t+1$; no time-uniform runtime follows.

For $\phi(s)=e^{-1/[s(1-s)]}$ on $(0,1)$, zero outside,
$N_\phi=\int_0^1\phi>2^{-11}$. Finite differentiation partitions
give, for $r\ge1$,
\begin{equation}
 \|\phi^{(r)}\|\le r!2^r(2r)^{2r}\le64^r(r!)^3,
 \qquad\|b_*^{(r)}\|\le2^{11}\|\phi^{(r-1)}\|.
 \label{eq:ev-flat-jets}
\end{equation}
Indeed $v=1/s+1/(1-s)$ has
$|v^{(j)}|\le j!v^{j+1}$ and differentiation bounds each expression
by $r!2^rv^{2r}e^{-v}$. The physical switch uses the separate
quotient of $e^{-1/s}$ and $e^{-1/(1-s)}$, with denominator
$\ge e^{-2}$; the same argument applies. Their efficient integrals
are paid below. Normalized values and jets retain the separate
$N_\phi$ floor and arithmetic debits.

To avoid undefined inactive components, use rational overlapping guards
inside the flat-zero cutoff region and above the saddle's exterior
onset. Skip $P$ on the inner side and $N$ near $t=0$; after $T_+$
the full-height late branch is paid. Harmless evaluations of $G$ on a
guard overlap enlarge its finite time/slack box and all evaluator caps.
This extends an exact-$G$ evaluation domain, not its comparison to $H$.

\subsection{The actual tuple and an explicit arithmetic allocation}
Actual $F$ spatial jets through four with one time derivative and $r$
spatial jets through three with one time derivative suffice; add one
time jet for coordinate moduli of a time row. $J_e=8$ covers the
source leaves. Every spatial, moving and temporal contribution survives.
For one blend $F=(1-h)F_0+hF_1$, for example,
\begin{align}
 B&=(1-h)B_0+hB_1+\tfrac12h_x(F_1-F_0),\notag\\
 E_{\rm comp}&=(1-h)E_0+hE_1,&
 E_{\rm temporal}&=h_t(F_1-F_0),\notag\\
 E_{\rm spatial}&=h_x(F_1-F_0)_z+\tfrac14h_{xx}(F_1-F_0).
 \label{eq:ev-defects}
\end{align}
Their sum is $F_t+F_{zz}$; $h_t$ includes the moving-gap term.
Apply the same identities recursively to nested blends.

For support plus coordinate slack put
\begin{align}
 X_q&=|x|+3,&T_q&=|t|+2,&B_q&=\max(97,\lceil T_q\rceil),\notag\\
 K_q&=8B_q(B_q+1)+36,&
 e_b&=4\eta e^{-4}(1+T_q^2)^{-1/4},\notag\\
 e_+&=e^{-3}K_q^{-3/2}e^{-X_q^2/72},&
 d_q&=\min(1,e_b,e_+)^2.
 \label{eq:ev-local-floor}
\end{align}
The actual regulator satisfies $q\ge d_q$ there. Geometric
interpolation is at least the smaller arm; the preceding source saddle
and complete $\Phi(0)$ bounds pay both arms. For $t\le64$ the
stronger time-only floor $\min(1,e_b)^2$ may be used. This is a local
evaluation floor; it does not replace the improved source-comparison
conditioning or assert a global floor after the switch.

There is a mechanical bound for every real gate. Give each input jet
caps $(V,L,A)$ for magnitude, first coordinate derivatives and error
amplification, initially $A=1$. Propagate
\begin{center}\small
\begin{tabular}{@{}llll@{}}\toprule
Gate & $V$ & $L$ & $A$\\\midrule
$u+v$ & $V_u+V_v$ & $L_u+L_v$ & $A_u+A_v+1$\\
$uv$ & $V_uV_v$ & $V_uL_v+V_vL_u$
 & $(V_u+1)A_v+(V_v+1)A_u+1$\\
$1/u$, $u\ge d$ & $1/d$ & $L_u/d^2$ & $2A_u/d^2+1$\\
$\log u$, $u\ge d$ & $V_u+1/d$ & $L_u/d$ & $2A_u/d+1$\\
$e^u$ & $e^{V_u}$ & $e^{V_u}L_u$ & $e^{V_u+1}A_u+1$\\\bottomrule
\end{tabular}
\end{center}
Split complex operations into real ones and implement complex inverse
through the positive modulus square. Induction proves error $\le A\tau$
when every input and gate debit is at most $\tau$, intermediate errors
are at most one and inverse/log inputs stay within $d/2$. For maxima
$V_*,L_*,A_*$ over the finite circuits and minimum floor $d_*$,
the choice
\begin{equation}
 \tau=\min\{e/(64A_*),1/(2A_*),d_*/(4A_*)\},\qquad
 p_{\rm arith}=\lceil\log_2(1/\tau)\rceil+2
 \label{eq:ev-arithmetic}
\end{equation}
proves integrand error $\le e/64$ and preserves every floor.
Magnitude storage also pays $\lceil\log_2(2+V_*)\rceil$ integer
bits. Provider internal precision is paid by its own recursion at this
output tolerance; $p_{\rm arith}$ is not a unit-cost oracle.

The raw integrand inventory, with the inherited exact definitions, is
\begin{center}\small
\begin{tabular}{@{}>{\raggedright\arraybackslash}p{.66\textwidth}>{\raggedright\arraybackslash}p{.28\textwidth}@{}}\toprule
Quantity & Kernel orders\\\midrule
$S=\log D$, $w=q/D$ & through four\\
$\rho$, direct $\rho_t$, $S_t$, $w_t$ & through two\\
$j=qF_{zz}\overline{F_z}/(\pi D^2)$ & through three internally\\
$C_{\rm reg},T_C,(C_{\rm reg})_t,(T_C)_t$ & through two\\
$b_{\rm reg}w$, actual $\mathcal S$, $V_0$, $c_rw$, $\Xi$ & through two\\
$T_U=\Xi/2+2\Re Z$ & through four\\
$D_E,D_B,T_J$ and separate defect addends & through two; $T_J$ through three\\\bottomrule
\end{tabular}
\end{center}
Direct time arithmetic uses $D_t=2\Re(F_t\bar F)+q_t$,
$S_t=D_t/D$, $w_t=q_t/D-qD_t/D^2$, and, with $a=F_z$,
\begin{equation}
 \rho_t=\frac{q_t|a|^2+2q\Re(a_t\bar a)}{\pi D^2}
                         -2\rho D_t/D.
 \label{eq:ev-direct-rhot}
\end{equation}
Thus $R_t$ is evaluated directly rather than defined by its PDE.
For $r_y=0$, the complete curvature time row is
\begin{align*}
 C_{\rm reg}&=(2w_xr_x+wr_{xx}-\chi_r r_x^2)/(4\pi),\quad\chi_r=w(1-w),\\
 (C_{\rm reg})_t&=[2w_{xt}r_x+2w_xr_{xt}+w_tr_{xx}+wr_{xxt}
       -(1-2w)w_tr_x^2-2\chi_r r_xr_{xt}]/(4\pi),\\
 w_{xt}&=q_{xt}/D-(q_xD_t+q_tD_x+qD_{xt})/D^2+2qD_xD_t/D^3.
\end{align*}
Differentiate the exact $T_C$ formula by the product rule for its time
row. Both curvature derivatives remain in the density equation.

For covariance apply the signed-stress identity first. For $|\alpha|\le2$,
\begin{align}
 P_\alpha={}&4\Re\langle K_\alpha,V_0\rangle
  -\langle K_{\alpha+xx}-K_{\alpha+yy},S\rangle
  +\langle K_\alpha,c_rw\rangle+\langle K_\alpha,\Xi\rangle,\notag\\
 Q_\alpha={}&P_\alpha-\sum_{i=x,y}s_i
       \sum_{\beta\le\alpha}\binom\alpha\beta
       U_{\beta+i}U_{\alpha-\beta+i},\quad s_x=1,\ s_y=-1.
 \label{eq:ev-actual-covariance}
\end{align}
These are actual full filtered mean products. If every elementary
filter/mean error is at most $e\le1$ and mean derivatives through
three are at most $M_{\rm mean}\ge1$, the complete Leibniz sums
give $|\operatorname{error}Q_\alpha|\le(8+24M_{\rm mean})e$.
$\Delta Q$ adds the two second-order rows with its own allocation.
No component covariance or separately positive pole square substitutes
for this signed recombination.

Apply the same gate rule to the finite post-filter circuit, treating
each raw integral as an input of error $e$, cap $4V$ and derivative
cap $4L$. Its maxima $A_{\rm post},V_{\rm post},L_{\rm post}$
are fixed before choosing $e$. For final scalar allowance $\zeta$, take
\begin{equation}
 e=\min\{1,\zeta/(16A_{\rm post}),
                 \zeta/[2^{20}(1+M_{\rm mean})^2]\}.
 \label{eq:ev-post-allocation}
\end{equation}
Rounding each post-filter gate with debit $e$ then gives final error
$\le A_{\rm post}e\le\zeta/16$. This includes all equation sums,
heat-potential correction, current derivatives and declared traces.
It does not presume a scalar target pays every equation component.

\begin{theorem}[Complete finite evaluator]\label{thm:complete-evaluator}
From a certified fixed descriptor and rational query/slack box, every
retained scalar of the full tuple has a computable finite absolute-error
evaluation algorithm, including the continuous filter, direct time
rows, actual covariance and all joining defects. For scalar comparison
targets $\theta_r$, put $\theta=\min(1,\min_r\theta_r)$ and
\begin{equation}
 \zeta(t,x)=\frac{\theta}{4(1+t^2)(1+x^2)}.
 \label{eq:ev-endpoint-schedule}
\end{equation}
Adding evaluator error at most $\zeta$ to the inherited analytical
reserve $\theta_r/2$ gives total $\le3\theta_r/4$, with separate
component allocations. Evaluation errors vanish at both infinite-time
and height ends. Boxes may be placed at every real height and time;
no uniform finite work budget over those boxes is asserted.
\end{theorem}
\begin{proof}
The source provider was proved above. One finite fallback gamma
provider on $\Re v>0$, $|v|\le A_v$, uses
\[
 L_N(v)=-\Log v-\gamma_Nv+
       \sum_{k=1}^N[v/k-\Log(1+v/k)],\quad
 \gamma_N=\sum_{k=1}^N1/k-\log N.
\]
For $N\ge2A_v$, its log-gamma value error is at most
$(A_v^2+A_v)/N$, first derivative error $(A_v+1)/N$, and
order-$r\ge2$ error $(r-2)!N^{1-r}$. The convergent product,
$|\gamma-\gamma_N|\le1/N$, $|z-\Log(1+z)|\le|z|^2$ and
integrated derivative tails prove these bounds \cite{DLMFEulerGamma}.
For example $N=\lceil2A_v+4(r_0!+1)(1+A_v)^2/\delta\rceil$
pays any fixed jet order $r_0$ with separate arithmetic reserve.
Elementary logarithms use their finite atanh series with the declared
right-half-plane floor; exponentials use factorial series, roots use
bisection or the right-half-plane log. The preceding saddle, phase
and profile recipes then terminate. Midpoint integration of $\phi$
has error at most $2/M$; separately approximate $N_\phi$ and the
partial integral, with ratio error at most $2^{13}\eta$ when their
errors are $\eta\le2^{-13}$. Tiny endpoint tails are bounded directly
by the exponential derivative estimate rather than inverted distances.

For a finite continuous-filter fallback, use interior polynomial
derivatives of $k(v)=(315/256)(1-v^2)^4$ on $[-1,1]$. Their caps
are
\[
 \kappa_i=\frac{315}{256}\sum_{2j\ge i}
               \binom4j(2j)!/(2j-i)!,\quad0\le i\le5.
\]
The fifth is only an interior quadrature derivative; the weak tuple
uses the zero extension only through order four. The gate rules give
a common spatial first-derivative cap $L_*$ for every raw integrand.
Tensor midpoint with $M_q=\max(1,\lceil64L_*/e\rceil)$ cells per
axis costs at most $4L_*/M_q\le e/16$, including variation between
nodes and across smooth seams. Node errors $e/64$ cost at most
$e/16$. Sum common-grid numerators exactly, multiply by $4/M_q^2$,
and round the final division within $e/8$. Thus each continuous raw
filter has error $<e$, and \eqref{eq:ev-post-allocation} finishes.

Every cap is a positive finite expression, not an unknown compact
supremum. For $G$ use \eqref{eq:ev-source-caps}; for the saddle use
its rectangle cap and a coarse gamma bound at tolerance $1/4$, then
positive finite jet recurrences. $N,q,r$ are explicit Gaussian or
polynomial leaves, and all profiles use \eqref{eq:ev-flat-jets}.
The inverse floors are \eqref{eq:ev-local-floor} or the stronger
paid branch floors. The provider hierarchy and circuit recursion are
finite and acyclic. Its direct finite cost is $M_q^2$ times provider
and scalar-gate costs plus exact-sum/post-filter costs; no special
function is counted as unit cost. This fallback may be exponential
in accuracy bits. The next subsection improves that dependence.

For supplied coordinate balls, propagate the first-derivative caps
through the same post-filter circuit to obtain $L_{\rm final}$. At
rational center $(t_0,x_0,y_0)$ set
$\zeta_0=\theta/[4(1+(|t_0|+1)^2)(1+(|x_0|+1)^2)]$ and require
each radius at most $\min(1,\zeta_0/[64(1+L_{\rm final})])$.
The coordinate contribution $\le3\zeta_0/64$, plus arithmetic
$\le\zeta_0/16$, stays below $\zeta_0$. Extra time jets supply
moduli for time rows. This is a contract on encoded coordinates.
Finally \eqref{eq:ev-endpoint-schedule} proves the stated limits.
\end{proof}

\subsection{Polynomial additional-accuracy work on a fixed box}
The preceding fallback is deliberately inefficient. Fix its descriptor
and compact rational query/slack box $\mathcal Q$ before requesting
additional accuracy. Write $\mathcal G_3(V,A)$ for pure-axis bounds
\begin{equation}
 |\partial_i^rf|\le VA^r(r!)^3\quad(r\ge0),\qquad V,A\ge1.
 \label{eq:pc-gevrey}
\end{equation}
Only pure $x$ and pure $y$ derivatives are needed for quadrature.
Each required fixed mixed spatial/time jet is first supplied as its
own leaf. Pure-axis bounds alone would not bound transverse derivatives.

The following closure rules prove explicit $V,A$ for all raw circuits:
addition uses $(V_f+V_g,\max A)$; product uses
$(3V_fV_g,\max A)$; one derivative uses $(V_fA_f,8A_f)$;
inverse with $f\ge d$ uses $(1/d,2A_f(1+V_f/d))$;
exponential uses $(e^{V_f},16A_f(1+V_f))$; logarithm uses
\[
 \left(\max(V_f+1/d,3V_fA_f/d),
       \max(8A_f,2A_f(1+V_f/d))\right).
\]
Pad all constants to one. For products,
$\sum_j\binom rj^{-2}\le3$ bounds the normalized Leibniz sum;
for derivatives use $(r+1)^3\le8^r$ for $r\ge1$. The inverse
recurrence is controlled by the geometric sum of
$(A_f/A_{\rm out})^j$, times $V_f/d\le A_{\rm out}/(2A_f)$.
The differential equation $(e^f)'=f'e^f$ gives the exponential rule;
apply $f'/f$ for the logarithm rule. These prove the claimed closure,
including $r=0$ by the value caps.

The leaves satisfy \eqref{eq:pc-gevrey}: $G$ follows by direct source
integration, with pure derivative growth at most $\max(1,U)^r$;
$N,q,r$ use their explicit formulas. Full-buffer analytic saddle
disks of radius $1/64$ and the explicit coarse arm/jet cap $M$ give
$M64^rr!$ by Cauchy, hence the weaker cubic-factorial bound. The
full late strip includes the glued real-height branch; its determinant
floor and holomorphy supply disks even when a real saddle lies on the
imaginary edge of its half-rectangle. Profiles satisfy
\eqref{eq:ev-flat-jets}; translating preserves their bounds and a
scale $h\le1$ replaces $A$ by $A/h$. Flat guarded extensions and
the cutoffs involving $|x|$ are constant near zero and smooth at joins.
Time rows retain $X'$ and all moving terms as fixed coefficients.
Thus the closure rules give common computable $V,A$ on the whole
support, independent of requested accuracy.

For an efficient gamma provider, request fixed derivatives through
$J_\Gamma$ to error $2^{-P}$ on $\Re z>1$ and put
\begin{align}
 m_\Gamma=N_\Gamma&=P+J_\Gamma+\lceil\log_2(J_\Gamma!)\rceil+8,\notag\\
 T_m(w)&=(w-1/2)\Log w-w+\tfrac12\log(2\pi)
       +\sum_{k=1}^{m-1}\frac{B_{2k}}{2k(2k-1)w^{2k-1}},\notag\\
 g(z)&=T_{m_\Gamma}(z+N_\Gamma)
                       -\sum_{a=0}^{N_\Gamma-1}\Log(z+a).
 \label{eq:pc-gamma}
\end{align}
Gamma recurrence preserves the normalized analytic logarithm exactly.
The sectorial Stirling bound and
$|B_{2m}|\le4(2m)!/6^{2m}$ give, on $|\zeta-z|\le1/2$,
$|R_m(\zeta+N)|\le(2/9)^m/m<2^{-2m}$; here
$\Re(\zeta+N)>m$ and $\sec^2(\arg(\zeta+N)/2)\le2$.
Cauchy then bounds derivative error by
$j!2^j2^{-2m}<2^{-P-4}$. This differentiates an analytic remainder,
not an asymptotic inequality \cite{DLMFGammaOuter2026,DLMFBernoulli}.

All elementary logarithms have geometric-precision algorithms. For
$w=x+iy$, $x>0$, binary-scale $x^2+y^2$ into $2^hs$, $1\le s<2$;
the atanh series for $\log s$ has ratio at most $1/3$. Invert
$|y|/x$ if needed, and for $0\le u\le1$ use
\[
 \arctan u=\arctan(1/2)
                 +\arctan((u-1/2)/(1+u/2)),
\]
whose Taylor arguments have modulus at most $1/2$. Machin $\pi$
and atanh $\log2$ were explicitly provided above. Bernoulli numbers
need no oracle: $b_n=(n+1)!B_n$ is the integer recurrence
$b_0=1$, $b_n=-\sum_{j<n}\binom{n+1}j[n!/(j+1)!]b_j$.
For input length $I_z$, $|z|\le A_z$, set
\[
 L=P+\lceil\log_2[32(A_z+2m+2)]\rceil+4,\qquad
 H_\Gamma=I_z+L+m+J_\Gamma+\lceil\log_2(A_z+m+4)\rceil+32.
\]
Round elementary logs to total complex error $2^{-L}$ on a common
grid and assemble the rational jet expressions exactly, then round
to $P+4$ fractional bits. The log coefficients sum in modulus to at
most $A_z+2m+2$ in the value row, at most one in the first derivative;
higher rows have rational powers only. Common denominators and factorial
lengths give sufficient bounds
$2^{100(J_\Gamma+1)}H_\Gamma^{10}$ work and
$2^{20(J_\Gamma+1)}H_\Gamma^3$ storage. Round nonrational saddle
inputs first, using the next fixed gamma derivative for displacement.
This is a log-gamma/jet bound. Absolute-error gamma values and
$\exp(\log S)$ retain their magnitude and cancellation payments.

For continuous filters, let $t_j=-1+2j/n$ and $\ell_j$ be their
Lagrange polynomials, $0\le j\le n$, $n\ge1$. Then
\[
 \sum_j|\ell_j(t)|\le2^nn^n/n!<6^n\quad(-1\le t\le1).
\]
Use exact rational weights $w_j=\int_{-1}^1\ell_j$ on $M$ panels
per axis. Taylor degree $n$ on a panel has remainder
$V(A/M)^{n+1}((n+1)!)^2$. Interpolation amplification and tensor
splitting give
\begin{equation}
 |I_xI_yf-Q_xQ_yf|
 \le4(1+6^n)^2V(A/M)^{n+1}((n+1)!)^2.
 \label{eq:pc-filter-remainder}
\end{equation}
For a raw target $e\le1$, choose
\begin{equation}
 n=\max(1,\lceil\log_2(64V/e)\rceil),\quad
 M=\max(1,\lceil2^{16}A(n+1)^2\rceil),\quad
 N_{\rm nodes}=[M(n+1)]^2\le2^{34}A^2(n+1)^6.
 \label{eq:pc-filter-nodes}
\end{equation}
Using $r!\le r^r$, the error is at most
$V2^{-10n-12}<e/16$. Only pure-axis derivatives were used.
The tensor weight norm is at most $4\cdot36^n$; hence node
error $e2^{-6n-10}$ contributes $<e/256$. Pass this accuracy through
\eqref{eq:ev-arithmetic}, preserving all magnitude and floor reserves.

Use the interior polynomial trace of each weak $K$ derivative at
endpoint nodes. Changing that boundary representative does not change
its continuous integral. High quadrature order differentiates this
interior polynomial integrand; it asserts no bounded fifth weak
derivative of the zero extension. Exact weights share denominator
$W=2^nn!(n+1)!$: write each numerator as
$\prod_{a\ne j}[nt-(2a-n)]$, with denominator
$2^n(-1)^{n-j}j!(n-j)!$, and integrate monomials. All monomial
denominators divide $(n+1)!$. Weight lengths are $O(n\log(n+1))$.
If node outputs share denominator $2^B$, the mapped tensor sum shares
\begin{equation}
 2^BM^2W^2.
 \label{eq:pc-sum-denominator}
\end{equation}
Stream integer numerators, rather than multiply a denominator per node,
and round the final division within $e/16$. Storage is
$O(B+\log M+n\log(n+1)+\log(2+V)+\log N_{\rm nodes}+n)$,
in addition to provider workspace. Thus raw filter error is $<e$,
uniformly between nodes and across the actual smooth seams.

The one-dimensional rule with $V=1,A=64$ evaluates $N_\phi$ and
partial integrals to $2^{-P}$ using $O((P+1)^3)$ nodes. Their
normalizer and ratio errors retain $N_\phi>2^{-11}$. For profile
derivatives through fixed $J_e$, if $v=1/[s(1-s)]$ satisfies
\[
 v\ge\max\{16J_e^2,2[B+J_e+\lceil\log_2(J_e!)\rceil+2]\},
\]
return zero with the proved leaf error $2^{-B}$; on the complement,
$v=O(B)$ and Bell/rational factors have polynomial length. Partial
integral tails include the normalizer and near one use one minus the
reflected tail. Endpoints are exact. The separate physical switch has
the same mechanism and its own paid denominator floor. This prevents
inverse-distance work at extremely small encoded endpoint distances.

At fixed descriptor/box, the source evaluator has degree-eleven
sufficient accuracy work, the saddle has $O(P+1)$ contraction steps
using the gamma provider, and phase inversion has $O(P+1)$ such
saddle calls. The physical real saddle has fixed contraction factor.
Elementary bounded-argument exponentials/trigonometric functions have
$O(P+1)$ Taylor length. The source constant $\Phi(0)$ needs only
$P+20$ modes: for $r\ge4$ its summand is at most
$64r^4e^{-3r^2}<2^{11-2r^2}\le2^{-r-4}$, since
$r^4\le24e^r$. Its retained exponential arguments are $O(P^2)$,
so their factorial lengths remain polynomial. Independent input
rounding is separately paid. Consequently all point providers fit
\begin{equation}
 C_{\rm point}(P,I)\le C_{\mathcal Q}(1+P+I)^{20},\qquad
 S_{\rm point}(P,I)\le C_{\mathcal Q}(1+P+I)^8.
 \label{eq:pc-point-cost}
\end{equation}
The two nested linear iterations give degree twelve before remaining
arithmetic; the displayed degree twenty is deliberately generous.

\begin{theorem}[Polynomial additional-accuracy cost at fixed data]
\label{thm:polynomial-accuracy}
Fix a certified descriptor and compact rational query/slack box
$\mathcal Q$. For rational/dyadic independent scalar parameters and
the explicit canonical providers, additional absolute accuracy $p$
bits and query encoding length $I$ have sufficient bounds
\begin{equation}
 C\le C_{\rm desc,\mathcal Q}(1+p+I)^{40},\qquad
 S\le C_{\rm desc,\mathcal Q}(1+p+I)^8.
 \label{eq:pc-complete-cost}
\end{equation}
For other computable fixed scalar inputs, let $C_{\rm ext}(P)$,
$S_{\rm ext}(P)$ bound one request to all their providers. Multiply
the work bound by
$1+C_{\rm ext}(C_{\rm desc,\mathcal Q}(1+p+I))$ and add
$S_{\rm ext}(C_{\rm desc,\mathcal Q}(1+p+I))$ to storage.
No uniform polynomial in height, time or descriptor length is asserted.
\end{theorem}
\begin{proof}
Use every circuit and actual mean product in
\eqref{eq:ev-actual-covariance}, with all direct time, curvature and
defect rows. The fixed post-filter caps make
$\log_2(1/e)=p+O_{\rm desc,\mathcal Q}(1)$ in
\eqref{eq:ev-post-allocation}. Thus $n=O(p+1)$,
$M=O((p+1)^2)$ and $N_{\rm nodes}=O((p+1)^6)$.
Node accuracy and each provider accuracy require $O(p+1)$ bits.
Node encoding length is $I+O(\log M+\log(n+1))$. Multiplying
the degree-twenty point cost by degree-six node count gives degree
twenty-six; degree forty pays weights, reading, exact summation,
rounding and post-filter gates with slack. Streaming nodes preserves
degree-eight storage. External scalar costs are explicitly substituted,
since a computable fixed scalar may have arbitrarily slow digits.
Coordinate balls require $p+O_{\rm desc,\mathcal Q}(1)$ bits by
the preceding finite derivative recursion, with no change of scope.
The endpoint schedule and total $3\theta_r/4$ comparison are unchanged.
\end{proof}

\subsection{A proved puncture bound and the remaining height consumer}
The inverse gate in the Gevrey rules still contains $V/d_q$. Hence
$A$ and the filter panel count can be exponential in height squared.
This is a sufficient cost, not a complexity lower bound. The local
model $F=a(z-c)$ with nearly constant $q$ has
$\rho(c)=|a|^2/(\pi q)=1/(\pi\delta^2)$,
$\delta=\sqrt q/|a|$: its narrow peak has finite mass.
It cannot be discarded to improve a quadrature estimate.

On holomorphic regions define smooth fractions
\[
 V_m=w^mF_{zz}\bar F/D,\qquad Z_m=w^mE\bar F/D,
 \qquad D=|F|^2+q,
\]
including at zeros. The exact identities are
\[
 S_t=-2\Re V_0+r_tw+2\Re Z_0,\quad
 w_t=r_tw(1-w)+2\Re V_1-2\Re Z_1,
\quad (w(1-w))_t=(1-2w)w_t.
\]
They retain actual $E$ and its separate addends.
\begin{lemma}[Weak-primitive punctures]\label{lemma:weak-punctures}
In a paid simple-root collar suppose $F=(z-a)b$, $|b|\ge s>0$,
$|F_{zz}/F|\le V_*/|z-a|+G_0$,
$|E/F|\le E_*/|z-a|+G_E$, and
$q_{\min}\le q\le q_{\max}$, $|F|\le F_{\max}$. Put
$B_{\log}=\max(1,|\log q_{\min}|,
|\log(F_{\max}^2+q_{\max})|)$. For a disk of radius $h$ in that
collar and any $m\ge0$,
\begin{align}
 \int_{|z-a|<h}|S|&\le\pi h^2B_{\log},\notag\\
 \int|V_m|&\le2\pi V_*h+\pi G_0h^2,&
 \int|Z_m|&\le2\pi E_*h+\pi G_Eh^2.
 \label{eq:pc-punctures}
\end{align}
These bounds have no inverse $q_{\min}$ factor.
\end{lemma}
\begin{proof}
The logarithm lies between the two stated endpoint logarithms.
Also $(1-w)w^m\le1$; multiply the pole-plus-regular estimates by
this factor and integrate $1/r$ in polar coordinates. Bounded
polynomials in $w,w(1-w)$ cost their bound times $\pi h^2$.
Multiplying by the actual weak test supremum pays its disk contribution.
For $N$ collars and coefficient cap $C$, choosing
$h\le e/[CN(1+B_{\log})]$ within the collar radius pays the
total after enlarging $C$ for displaced-disk guards.
\end{proof}
The first unpaid height-cost theorem must consume this lemma for the
complete inventory: construct and guard every intersecting collar,
bound normalized retained providers away from punctures, retain all
joining-test commutators and nonholomorphic $B$ terms, and replace
pointwise post-filter caps by primitive $L^1$ and actual mean bounds.
Coordinate balls additionally need highest-kernel boundary traces and
time moduli: differentiating the fourth weak kernel derivative can
produce boundary line measures. The present lemma does not pay them.
First derive the exact weak representation; deleting disks from raw
$\rho,j,D_E$ would lose mass or signed dipoles. Neither this lemma
nor efficient gamma alone proves uniformly affordable all-height work.
